\chapter[APRESENTAÇÃO DA TRAJETÓRIA DO ALUNO]{APRESENTAÇÃO DA TRAJETÓRIA DO ALUNO}

Minha trajetória até a Engenharia de Software começou com uma mudança de planos. Em 2021, enquanto cursava o terceiro ano do ensino médio, minha intenção era seguir carreira na Medicina. Entretanto, naquele período, fiquei muito doente em decorrência de uma doença autoimune e precisei reconsiderar o caminho que havia imaginado para os anos seguintes. Esse momento trouxe a necessidade de olhar para minhas condições de saúde e pensar em uma formação que também fosse compatível com a realidade que eu passava a enfrentar.

A possibilidade de trabalhar de forma remota ou híbrida foi um dos fatores que me aproximaram da área de tecnologia. Como eu já passava bastante tempo utilizando computadores, a Engenharia de Software surgiu como uma alternativa que reunia familiaridade com o ambiente digital e perspectivas de maior flexibilidade profissional. Assim, ingressei no curso no segundo semestre de 2022. Apesar dessa proximidade com computadores, eu ainda não possuía experiência em desenvolvimento de software e comecei a graduação com muitos conhecimentos a construir.

Ao longo do curso, a AGES se tornou um espaço importante para transformar esse contato inicial em experiência prática. No projeto Plan Line, como AGES I, tive minha primeira vivência mais concreta com o desenvolvimento de uma aplicação em equipe. Foi nesse período que comecei a compreender as dificuldades de aprender tecnologias novas enquanto precisava contribuir com entregas reais. Também descobri maior afinidade com o backend, interesse que se aprofundou na AGES II, no projeto de Certificação por meio de Blockchain, especialmente por meio da modelagem de dados e do desenvolvimento com Java e Spring Boot.

Em 2024, comecei a procurar uma oportunidade de estágio e, no início do segundo semestre, ingressei na SAP, no time de Globalization. Minha atuação esteve relacionada à localização de um produto de folha de pagamento para o mercado espanhol. Essa foi minha primeira experiência profissional na área e me colocou em contato com um contexto de desenvolvimento bastante específico, no qual o software precisava atender às particularidades de um determinado mercado. A experiência ampliou minha percepção de que desenvolver sistemas também exige compreender o domínio em que eles serão utilizados.

Com o passar do tempo, percebi que a especialização daquele produto, desenvolvido em ABAP e em um ambiente on-premises, não correspondia ao direcionamento que eu desejava dar à minha formação. Aproximadamente um ano após minha entrada, decidi deixar a empresa. Essa decisão fez parte do processo de entender melhor meus interesses profissionais e de reconhecer que as escolhas de carreira precisam considerar tanto as oportunidades disponíveis quanto o tipo de experiência que se deseja construir.

Em fevereiro de 2026, iniciei uma nova experiência em uma empresa local de desenvolvimento de hardware que buscava ingressar no mercado de software, onde continuo trabalhando. A mudança me aproximou de um contexto empresarial diferente daquele encontrado na SAP e passou a compor uma nova etapa da minha trajetória profissional. Ao conciliar esse trabalho com a graduação, pude continuar aproximando os conhecimentos acadêmicos das situações encontradas no mercado.

Nas passagens seguintes pela AGES, minha participação também se transformou. No VínculoHub Portal, como AGES III, passei a contribuir com mais segurança técnica e a perceber com maior clareza a importância do interesse pelo negócio, da comunicação e da organização das equipes. Já no ReadRace, como AGES IV, assumi responsabilidades de gestão, acompanhamento de tarefas e apresentações aos stakeholders, além de continuar contribuindo com desenvolvimento e revisão de código. Essa atuação trouxe desafios diferentes dos iniciais, especialmente na necessidade de orientar colegas e comunicar dificuldades antes que comprometessem as entregas.

Ao olhar para esse percurso, percebo uma formação construída por adaptações, aprendizado e escolhas progressivamente mais conscientes. A Engenharia de Software entrou na minha vida em um momento em que precisei rever meus planos, mas, ao longo das experiências acadêmicas e profissionais, fui encontrando interesses próprios dentro da área. As quatro passagens pela AGES fazem parte desse amadurecimento: do primeiro contato com um projeto real à responsabilidade de acompanhar uma equipe, cada etapa contribuiu para ampliar minha compreensão sobre o trabalho e sobre o profissional que desejo me tornar.
    
